\section{Results, Limitations, and Further Opportunities}
\label{sec:results}

\subsection{What We Validated}

\begin{enumerate}
\item \textbf{The concept is technically viable.} The KGDRL architecture successfully generated feasible scheduling recommendations and supported scenario-based decision making, demonstrating that AI-assisted wave allocation can be implemented within a lightweight prototype.

\item \textbf{Reducing perceived risk matters more than improving performance.} In informal testing with supply-chain classmates, the What-If simulator generated significantly more interest than benchmark results alone. Users were more willing to explore the system when they could safely experiment with scenarios before committing to adoption.

\end{enumerate}

\subsection{An Unexpected Finding}

One unexpected lesson was that increasing algorithmic sophistication did not automatically increase perceived value. Discussions during development and Demo Day focused far more on deployment effort, explainability, and integration concerns than on optimization performance. This reinforced our central insight that organizational attention, rather than algorithm quality, represents the primary constraint on adoption.

\subsection{Current Limitations}

\begin{enumerate}
\item \textbf{Prototype-level infrastructure.} Streamlit enabled rapid development but lacks the responsiveness and scalability expected in a production warehouse environment.

\item \textbf{Simplified operational assumptions.} The reward function captures key scheduling objectives but omits real-world factors such as picker fatigue, equipment congestion, and organizational coordination.

\item \textbf{Limited interactivity.} Several modules remain demonstration-oriented rather than fully interactive, restricting their usefulness for continuous decision support.


\end{enumerate}

\subsection{Next Step}

The next step is not further algorithm validation, but adoption validation. A four-week pilot in a real warehouse would allow observation of whether managers regularly use the system, trust its recommendations, and integrate it into daily decision-making. Such evidence would directly test whether organizational attention is indeed the primary barrier to adoption.

The central lesson of Maverick-SORT is that creating a technically capable AI system is only part of the innovation challenge. Real impact depends on whether potential users can understand, evaluate, and trust that system within the constraints of their everyday work.